\DocumentMetadata{
  pdfversion=2.0,
  pdfstandard=ua-2,
  lang=en-US,
  tagging=on,
  tagging-setup={math/setup=mathml-SE,tabsorder=structure}
}
\documentclass[11pt,letterpaper]{article}
\usepackage[margin=0.68in,includefoot]{geometry}
\usepackage{newcomputermodern}
\usepackage{amsmath,amscd,mathtools}
\usepackage{tikz,adjustbox}
\providecommand{\square}{\mdlgwhtsquare}
\usepackage[hidelinks]{hyperref}
\hypersetup{
  pdftitle={Analysis PhD Exam, May 1997},
  pdfauthor={University of Florida Department of Mathematics},
  pdfsubject={Graduate mathematics examination},
  pdfdisplaydoctitle=true,
  bookmarksopen=true
}
\setcounter{secnumdepth}{0}
\setlength{\parindent}{0pt}
\setlength{\parskip}{0.28em}
\setlength{\emergencystretch}{3em}
\setlength{\leftmargini}{2.15em}
\setlength{\leftmarginii}{2.65em}
\setlength{\leftmarginiii}{3em}
\pagestyle{empty}
\raggedbottom
\allowdisplaybreaks
\NewTaggingSocketPlug{block/list/label}{examproblem}{%
  \tagstructbegin{tag=Lbl}\tagstructend%
  \tagstructbegin{tag=\LBody}%
  \tagstructbegin{tag=\ProblemHeadingTag,title-o={\ProblemHeadingText}}%
  \tagmcbegin{tag=\ProblemHeadingTag,actualtext={\ProblemHeadingText}}%
  #2%
  \tagmcend%
  \tagstructend%
}
\newenvironment{examproblems}[2]{%
  \def\ProblemHeadingTag{#2}%
  \AssignTaggingSocketPlug{block/list/label}{examproblem}%
  \begin{enumerate}%
  \setcounter{enumi}{#1}%
  \renewcommand{\labelenumi}{\textbf{\theenumi.}}%
  \setlength{\topsep}{0.35em}%
  \setlength{\partopsep}{0pt}%
  \setlength{\itemsep}{0.48em}%
  \setlength{\parsep}{0.18em}%
}{\end{enumerate}}
\newenvironment{examsubparts}{%
  \AssignTaggingSocketPlug{block/list/label}{default}%
  \begin{enumerate}%
  \setlength{\topsep}{0.18em}%
  \setlength{\partopsep}{0pt}%
  \setlength{\itemsep}{0.16em}%
  \setlength{\parsep}{0.1em}%
}{\end{enumerate}}
\newenvironment{examinstructions}{%
  \par\begingroup\itshape
}{%
  \par\endgroup\vspace{0.18em}
}
\makeatletter
\renewcommand\subsection{\@startsection{subsection}{2}{0pt}%
  {-1.25ex plus -.25ex minus -.1ex}%
  {0.55ex plus .1ex}%
  {\normalfont\large\bfseries}}
\renewcommand{\@maketitle}{%
  \newpage\null\vskip -1em%
  {\footnotesize\bfseries
    \href{https://www.ufl.edu/}{University of Florida}, \href{https://math.ufl.edu/}{Department of Mathematics}\par}%
  \vskip -0.5em%
  \tagstructbegin{tag=H1,title={Analysis PhD Exam, May 1997}}%
  \tagpdfparaOff%
  \tagmcbegin{tag=H1}%
    {\LARGE\bfseries\href{https://gma.math.ufl.edu/exams/analysis-phd/}{Analysis PhD Exam}, May 1997\par}%
  \tagmcend%
  \tagpdfparaOn%
  \tagstructend%
  \vskip 0.25em%
  \tagmcbegin{artifact}%
    \hrule height 1pt%
  \tagmcend%
  \vskip 1em%
}
\makeatother
\title{Analysis PhD Exam, May 1997}
\author{}
\date{}
\begin{document}
\maketitle
\thispagestyle{empty}
\begin{examinstructions}
Be sure to write all steps in a neat and logical fashion to receive credit. Give reasons for your steps.

\par
Definition: A σ-algebra \(\mathcal F\) is separable if \(\mathcal F=\sigma\{A_j:j\geq 1\}\) for some sequence of sets \(A_j\).
\end{examinstructions}
\begin{examproblems}{0}{H2}
\def\ProblemHeadingText{Problem 1.}
\item
Let \(\{\mathcal F_k\}\) be a sequence of separable σ-algebras of subsets of a set~\(X\). Let \(\mathcal F\) be the smallest σ-algebra of subsets of~\(X\) which contains \(\mathcal F_k\) for each~\(k\). Prove that \(\mathcal F\) is separable. [Hint: Use the monotone class theorem.]
\def\ProblemHeadingText{Problem 2.}
\item
Let \((X,\mathcal F,\mu)\) be a finite measure space and let \((Y,\Sigma)\) be a measurable space. Let \(f:X\to Y\) be \(\mathcal F,\Sigma\)-measurable. Define \(\nu(B)=\mu(f^{-1}(B))\) for \(B\in\Sigma\). Show~\(\nu\) is a measure on~\(\Sigma\), and if \(g:Y\to[0,\infty]\) is~\(\Sigma\)-measurable, then \(\int_Y g\,d\nu=\int_X g\circ f\,d\mu\).
\def\ProblemHeadingText{Problem 3.}
\item
Let~\(m\) be Lebesgue measure on \(\mathbb R\). Let \(f_m=-\mathbf 1_{A_m}\), where \(A_m=(m,\infty)\). Show \((f_m)\) is increasing, but the monotone convergence theorem is not true. Why is this?
\def\ProblemHeadingText{Problem 4.}
\item
Let \((X,\mathcal S,\mu)\) and \((Y,\mathcal J,\nu)\) be two finite measure spaces. Construct the product measure \(\mu\times\nu\) on \(\mathcal S\otimes\mathcal J\). Be sure to prove \(\mu\times\nu\) is countably additive.
\def\ProblemHeadingText{Problem 5.}
\item
Let \((X,\Sigma,\mu)\) be a finite measure space. Prove that the dual of \(L^1(X,\Sigma,\mu)\) is isometrically isomorphic to \(L^\infty(X,\Sigma,\mu)\).
\def\ProblemHeadingText{Problem 6.}
\item
Let \((X,\mathcal F,\mu)\) be a finite measure space. Suppose that \(\mathcal F\) is separable. Show \(L^1(X,\mathcal F,\mu)\) is separable.
\def\ProblemHeadingText{Problem 7.}
\item
Let \((X,\mathcal F,\mu)\) be a measure space. Let \(\{f_m\}\) be a sequence of measurable functions such that \(\int_X|f_m|\,d\mu<\frac{1}{2^m}\). Does it follow that \(\sum_{m=1}^\infty f_m\) converges a.e. \(\mu\) on~\(X\)?
\def\ProblemHeadingText{Problem 8.}
\item
Let \((\Omega,\mathcal F,P)\) be a probability space. Let \(X:\Omega\to\mathbb R\) be an integrable function. Let \(F(x)=P(X<x)\), for \(x\in\mathbb R\).
\begin{examsubparts}
\item[\textnormal{(a)}]
Show \(F:\mathbb R\to[0,1]\) is monotone, left continuous, \(\lim_{x\to-\infty}F(x)=0\), and \(\lim_{x\to\infty}F(x)=1\).
\item[\textnormal{(b)}]
Show \(\int_\Omega X^2\,dP=\int_{\mathbb{R}} x^2\,dF(x)\).
\end{examsubparts}
\def\ProblemHeadingText{Problem 9.}
\item
Let \((\Omega,\mathcal F,P)\) be a probability space. Let \((X_m)\) be a sequence of independent random variables on~\(\Omega\).
\begin{examsubparts}
\item[\textnormal{(a)}]
Show \((\overline{\lim}X_m>3)\) is a tail event relative to \((X_m)\).
\item[\textnormal{(b)}]
If \(P(\overline{\lim}X_m>3)>0.2\), what is \(P(\overline{\lim}X_m>3)\)?
\end{examsubparts}
\def\ProblemHeadingText{Problem 10.}
\item
Let \(\ell^\infty\) be the Banach space consisting of all real sequences \((a_i)\) such that \(\sup_i|a_i|<\infty\), with \(\|(a_i)\|:=\sup_i|a_i|\). Let~\(c\) be the subspace of \(\ell^\infty\) consisting of \((a_i)\) such that \(\lim a_i\) exists. Prove that there exists a continuous linear functional \(F:\ell^\infty\to\mathbb R\) such that \(F((a_i))=\lim a_i\) for each \((a_i)\in c\). [Hint: Use the Hahn–Banach theorem.]
\end{examproblems}
\end{document}
